\chapter{Domain-Specific FRFP Implementation Playbooks}
\label{ch:frfp-domain-playbooks}

This chapter gives concrete, domain-specific playbooks for implementing
FRFP in three high-impact areas:

\begin{enumerate}
  \item scientific and engineering research;
  \item finance and risk management;
  \item clinical practice and healthcare delivery.
\end{enumerate}

Each playbook is structured as a short, practical guide:

\begin{itemize}
  \item typical explicit and tacit spaces in the domain;
  \item recommended boundary design;
  \item safe-horizon choices;
  \item governance patterns;
  \item a compact implementation checklist.
\end{itemize}

The aim is not to exhaust domain detail, but to show how the general FRFP
architecture can be concretely instantiated.

\section{Playbook I: Scientific and Engineering Research}
\label{sec:playbook-research}

\subsection{Typical Explicit and Tacit Spaces}

\paragraph{Explicit space \(E_{\mathrm{res}}\).}

\begin{itemize}
  \item datasets, experimental logs, simulation outputs;
  \item code, models, pipelines, configuration files;
  \item manuscripts, preprints, lab notebooks, design docs;
  \item metrics and benchmarks (loss curves, accuracy, error bars).
\end{itemize}

\paragraph{Tacit space \(T_{\mathrm{res}}\).}

\begin{itemize}
  \item researchers’ intuitions about what questions are important or
        tractable;
  \item judgement about experimental design, artefact quality, and
        whether results ``make sense'';
  \item field-wide norms about evidence, novelty, and acceptable risk;
  \item institutional incentives (career, funding, publication venues).
\end{itemize}

\subsection{FRFP Boundary Design in Research}

\paragraph{Key boundaries.}

\begin{itemize}
  \item \textbf{Research claims}:
        when a result moves from ``internal exploration'' to
        ``external claim'' (preprint, paper, patent).
  \item \textbf{Code and model releases}:
        when an internal tool becomes a public artefact or shared
        infrastructure.
  \item \textbf{Project commitments}:
        when a lab or organisation commits resources to a long-horizon
        programme (e.g.\ a multi-year line of work or experiment).
\end{itemize}

\paragraph{Design moves.}

\begin{itemize}
  \item Require that any \emph{claim-bearing document} (paper, report,
        safety case) be explicitly marked as such and go through a
        human review pipeline distinct from AI-assisted drafting.
  \item Maintain explicit provenance:
        AI-generated analyses and text are tagged, so that reviewers
        know when they are endorsing machine-produced content.
  \item For code/model releases, implement a release gate where
        technical leads and domain experts jointly sign off on:
        \begin{itemize}
          \item domain of validity;
          \item known failure modes;
          \item recommended and prohibited uses.
        \end{itemize}
\end{itemize}

\subsection{Safe Horizons in Research Programmes}

\paragraph{Temporal horizons.}

\begin{itemize}
  \item Set explicit review intervals for large programmes:
        e.g.\ every 6--12 months, with an external or cross-team review
        of assumptions, goals, and methods.
  \item For high-risk directions (e.g.\ dual-use capabilities),
        enforce shorter horizons and require renewal of approval.
\end{itemize}

\paragraph{Scale horizons.}

\begin{itemize}
  \item Limit the number of downstream projects that can depend on a
        single unvalidated model or dataset before more thorough
        validation is performed.
  \item For AI tools used across a lab or organisation, cap adoption
        (e.g.\ number of users or critical workflows) until an initial
        evaluation phase is complete.
\end{itemize}

\paragraph{Domain horizons.}

\begin{itemize}
  \item Require explicit justification before using a tool outside the
        domain where it was validated (e.g.\ reusing a model trained
        for one regime in a qualitatively different one).
  \item Trigger re-grounding when key assumptions (hardware, scale,
        data sources, institutional context) change.
\end{itemize}

\subsection{Governance Patterns for Research}

\begin{itemize}
  \item Establish an internal \textbf{FRFP board} (or adapt an existing
        research committee) with responsibility for:
        \begin{itemize}
          \item approving high-impact projects and releases;
          \item setting safe-horizon policies;
          \item reviewing incident reports and near-misses.
        \end{itemize}
  \item Integrate FRFP concepts into:
        \begin{itemize}
          \item lab meeting culture (explicit discussion of boundaries
                and tacit judgement);
          \item mentoring and training (how to use AI tools without
                eroding core skills);
          \item collaboration agreements (roles and responsibilities for
                human and AI components).
        \end{itemize}
  \item Align with external institutions:
        IRBs, safety review boards, professional societies and journals
        (e.g.\ disclosure of AI-generated analyses and code).
\end{itemize}

\subsection{Implementation Checklist (Research)}

For a research group or engineering team:

\begin{enumerate}
  \item \textbf{Write down \(E_{\mathrm{res}}\) and \(T_{\mathrm{res}}\)} for your
        main projects.
  \item \textbf{Mark boundaries}:
        which artefacts count as claims, releases, or commitments?
  \item \textbf{Define safe horizons} for:
        \begin{itemize}
          \item programme duration before review;
          \item scale of deployment of internal tools;
          \item reuse of models outside their training domain.
        \end{itemize}
  \item \textbf{Set up a lightweight FRFP board} (even 2--3 people) to
        approve boundary events and horizon extensions.
  \item \textbf{Instrument provenance}:
        tag AI-generated analyses and text; log endorsements at
        boundaries.
\end{enumerate}

\section{Playbook II: Finance and Risk Management}
\label{sec:playbook-finance}

\subsection{Typical Explicit and Tacit Spaces}

\paragraph{Explicit space \(E_{\mathrm{fin}}\).}

\begin{itemize}
  \item trades, orders, positions, and P\&L records;
  \item risk metrics:
        VaR, stress tests, greeks, liquidity measures;
  \item client data, credit histories, transaction logs;
  \item models and algorithms:
        pricing models, risk engines, execution algos, ML forecasters;
  \item policies and limit tables implemented in code or configuration.
\end{itemize}

\paragraph{Tacit space \(T_{\mathrm{fin}}\).}

\begin{itemize}
  \item traders’ and portfolio managers’ market intuition;
  \item risk officers’ judgement about model limitations and stress
        scenarios;
  \item executives’ and boards’ beliefs about risk appetite and
        strategy;
  \item regulators’ evolving expectations;
  \item client and market participants’ perceptions of stability and
        fairness.
\end{itemize}

\subsection{FRFP Boundary Design in Finance}

\paragraph{Key boundaries.}

\begin{itemize}
  \item \textbf{Capital allocation and limits}:
        approvals for large positions, new strategies, or structural
        hedges.
  \item \textbf{Model changes}:
        implementation or retirement of risk/pricing models.
  \item \textbf{Regulatory and public disclosures}:
        risk reports, investor communications.
\end{itemize}

\paragraph{Design moves.}

\begin{itemize}
  \item Treat \emph{model approval} as a major boundary event:
        \begin{itemize}
          \item independent model validation;
          \item parallel runs with legacy models;
          \item sign-off by risk and business leadership.
        \end{itemize}
  \item Make limit changes and waivers explicit boundary events:
        who approved, under what rationale, with what horizon.
  \item Distinguish between \emph{hedging} and \emph{directional}
        strategies in both code and governance, rather than relying on
        labels alone.
\end{itemize}

\subsection{Safe Horizons in Financial Systems}

\paragraph{Temporal horizons.}

\begin{itemize}
  \item Set review cycles for strategies (e.g.\ quarterly or annual
        deep dives) independent of short-term P\&L.
  \item For complex or illiquid portfolios, require more frequent
        stress tests and scenario reviews.
\end{itemize}

\paragraph{Scale horizons.}

\begin{itemize}
  \item Cap notional exposure and risk measures per strategy or book
        until performance and model fit are demonstrated.
  \item Tie horizon extensions to both quantitative evidence and tacit
        review of market structure and counterparty behaviour.
\end{itemize}

\paragraph{Domain horizons.}

\begin{itemize}
  \item Restrict model use to markets, instruments, and regimes where
        assumptions hold; treat entry into new asset classes or
        liquidity regimes as new domains requiring fresh validation.
  \item For ML models, monitor for data drift and structural breaks
        (e.g.\ regime shifts, new regulations); trigger re-grounding.
\end{itemize}

\subsection{Governance Patterns for Finance}

\begin{itemize}
  \item Define a \textbf{FRFP-aware risk committee} with authority over:
        \begin{itemize}
          \item model approvals and retirements;
          \item safe-horizon definitions for strategies;
          \item limit frameworks and escalation procedures.
        \end{itemize}
  \item Ensure clear separation of roles:
        \begin{itemize}
          \item traders and quants design strategies and models;
          \item independent risk validates and challenges;
          \item the committee arbitrates and sets boundaries.
        \end{itemize}
  \item Align with external regulators:
        treat regulatory stress tests and model-risk rules as
        institutional operators in FRFP terms, and design internal
        processes to be legible to them.
\end{itemize}

\subsection{Implementation Checklist (Finance)}

For a trading business, risk group, or fintech:

\begin{enumerate}
  \item \textbf{Catalogue key models and metrics} in
        \(E_{\mathrm{fin}}\), including their intended semantics and
        data sources.
  \item \textbf{Map tacit risk decisions} in \(T_{\mathrm{fin}}\):
        who actually decides on exposures, hedges, and new products.
  \item \textbf{Define boundaries}:
        \begin{itemize}
          \item model approval / retirement;
          \item strategy launch / scale-up;
          \item limit changes and waivers.
        \end{itemize}
  \item \textbf{Set safe horizons} for:
        \begin{itemize}
          \item maximum exposure pre-review;
          \item maximum time between independent stress-testing;
          \item valid domain of each model.
        \end{itemize}
  \item \textbf{Establish or adapt a risk committee} to act as a
        governance operator with real authority over the above.
\end{enumerate}

\section{Playbook III: Clinical Practice and Healthcare Delivery}
\label{sec:playbook-clinical}

\subsection{Typical Explicit and Tacit Spaces}

\paragraph{Explicit space \(E_{\mathrm{clin}}\).}

\begin{itemize}
  \item electronic health records (EHRs), lab results, imaging data;
  \item clinical guidelines and order sets;
  \item AI models for:
        \begin{itemize}
          \item diagnosis and screening (e.g.\ DR detection, imaging
                triage);
          \item risk prediction (readmission, deterioration, sepsis);
          \item workflow optimisation and resource allocation.
        \end{itemize}
  \item decision-support outputs:
        alerts, risk scores, recommendations, triage flags, auto-draft
        notes.
\end{itemize}

\paragraph{Tacit space \(T_{\mathrm{clin}}\).}

\begin{itemize}
  \item clinicians' tacit knowledge of patients, disease courses,
        local practice norms;
  \item nurses’ and allied staff’s understanding of workflows and
        practical constraints;
  \item institutional culture around safety, liability, and autonomy;
  \item patient trust, preferences, and understanding of care;
  \item regulator and payer expectations.
\end{itemize}

\subsection{FRFP Boundary Design in Clinical Settings}

\paragraph{Key boundaries.}

\begin{itemize}
  \item \textbf{Diagnostic and treatment decisions}:
        diagnoses entered, treatments ordered, prescriptions signed.
  \item \textbf{Screening and triage decisions}:
        who is referred, prioritised, or discharged.
  \item \textbf{Documentation and billing}:
        what goes into the official record and claims.
\end{itemize}

\paragraph{Design moves.}

\begin{itemize}
  \item Treat AI outputs as structured inputs to clinical reasoning,
        not as replacements for it:
        UIs should present recommendations alongside underlying data
        and uncertainty.
  \item For autonomous systems (like DR screening AI), make the AI’s
        \emph{scope} explicit:
        the specific question it answers, and what it does \emph{not}
        decide (e.g.\ treatment).
  \item Require explicit clinician endorsement at key boundaries:
        entering a diagnosis, confirming a treatment plan, signing a
        discharge.
\end{itemize}

\subsection{Safe Horizons in Clinical Programmes}

\paragraph{Temporal horizons.}

\begin{itemize}
  \item For screening and risk-prediction programmes, define
        re-evaluation intervals (e.g.\ annually) for both:
        \begin{itemize}
          \item model performance;
          \item workflow and outcome impact (over- and under-treatment).
        \end{itemize}
  \item For AI tools used at the bedside, schedule periodic
        clinician-survey cycles and qualitative feedback rounds.
\end{itemize}

\paragraph{Scale horizons.}

\begin{itemize}
  \item Roll out new AI tools in phases:
        pilot units $\to$ hospital-wide $\to$ system-wide, with clear
        criteria for progression.
  \item Limit their use in high-acuity or vulnerable populations until
        sufficient evidence is available.
\end{itemize}

\paragraph{Domain horizons.}

\begin{itemize}
  \item Do not deploy models outside the populations in which they were
        validated (e.g.\ different hospitals, countries, or demographic
        mixes) without adaptation and revalidation.
  \item Explicitly track changes in clinical practice patterns,
        technology (e.g.\ new imaging devices), and epidemiology that
        may invalidate prior validations.
\end{itemize}

\subsection{Governance Patterns for Clinical AI}

\begin{itemize}
  \item Use or extend existing structures:
        \begin{itemize}
          \item clinical governance and safety committees;
          \item quality-improvement (QI) and patient-safety teams;
          \item IRBs and ethics boards for research deployments.
        \end{itemize}
  \item Give these bodies explicit FRFP responsibilities:
        \begin{itemize}
          \item approve indications-for-use and deployment domains;
          \item set safe horizons and monitoring plans;
          \item review incidents and near-misses involving AI tools.
        \end{itemize}
  \item Align with external operators:
        regulators (e.g.\ device approvals), payers (reimbursement), and
        professional societies (guidelines) as external FRFP operators
        that can change admissible runs.
\end{itemize}

\subsection{Implementation Checklist (Clinical)}

For a hospital, clinic, or health system:

\begin{enumerate}
  \item \textbf{Inventory AI tools} in
        \(E_{\mathrm{clin}}\):
        what models exist, what they predict, and where they are used
        in workflows.
  \item \textbf{Map tacit actors} in
        \(T_{\mathrm{clin}}\):
        clinicians, nurses, administrators, patients, regulators.
  \item \textbf{Define boundaries}:
        \begin{itemize}
          \item which decisions AI may take autonomously (if any);
          \item which require clinician endorsement;
          \item what gets written into the medical record.
        \end{itemize}
  \item \textbf{Set safe horizons}:
        \begin{itemize}
          \item pilot sizes and durations;
          \item review intervals for performance and impact;
          \item domain limits (sites, populations, devices).
        \end{itemize}
  \item \textbf{Empower a clinical AI committee} to:
        \begin{itemize}
          \item approve tools and indications;
          \item oversee monitoring and incident review;
          \item communicate with external regulators and payers.
        \end{itemize}
\end{enumerate}

\section{Using the Playbooks}

These playbooks are deliberately schematic.
For each domain, they provide:

\begin{itemize}
  \item a first pass at identifying \(E\), \(T\), boundaries, horizons,
        and governance operators;
  \item concrete steps an organisation can take in months, not years,
        to bring existing or planned systems closer to FRFP alignment.
\end{itemize}

In practice, implementing FRFP will involve adapting these templates to
local context and constraints.
The case studies in Volume~II provide detailed illustrations of how
similar patterns have failed or succeeded in related systems.
By combining the playbooks with the pattern language in
Chapter~\ref{ch:frfp-patterns} and the methodology in
Chapter~\ref{ch:frfp-method}, practitioners can build domain-specific
FRFP architectures that are both principled and operational.
